\section{Conclusion}

In this work, we are the first to study how test-time compute approaches, particularly large reasoning models, handle contextual privacy in probing and agentic settings. Our experiments on a suite of 13 models reveal that, while reasoning traces are key to increasing capability, they pose a new and overlooked privacy risk. These traces are often rich in personal data and can easily leak into the final output, either accidentally or via prompt injection attacks. While increasing the test-time compute budget makes the model more private in its final answer, it enriches its easily accessible reasoning over sensitive data.
We argue that future research should prioritise mitigation and alignment strategies to protect both the reasoning process and the final outputs. This includes extending efforts like \citet{jiang2025safechain}, which focus on jailbreak attacks, to also address privacy concerns. Moreover, advances in efficient reasoning \citep{sui2025stopoverthinkingsurveyefficient} may help reduce the exposure risk by naturally limiting the length and verbosity of reasoning traces.